% TOP_PROBLEM: {"id": 30006070, "title": "The Generalized Étale-Homotopy Section Conjecture", "category_id": 5, "category": {"id": 5, "name": "algebraic_geometry", "display_name": "Algebraic Geometry", "description": "Geometric objects defined by polynomial equations.", "slug": "algebraic-geometry", "order_index": 5, "created_at": "2026-07-31T15:26:25.671Z"}, "status": "open"}
% FIRST_POSTED: 2026-10-06
% ATTEMPT_STATUS: unresolved
% OPENAI_CLAIM: {"schema_version":1,"source_repo":"https://github.com/openai/math","source_commit":"adc7f1241b42e322a6451854ab7e4b4c146bf78a","release_date":"2026-10-06","match":"partial","resolution":"partial","independently_reviewed":false,"second_pass":"n/a","families":[{"family":"019","title":"The local <i>p</i>-adic section conjecture and global consequences","subject":"Number theory","manuscripts":[{"title":"The p-adic section conjecture","pdf_path":"preprints/The-p-adic-section-conjecture-September-24-2026/main.pdf","readme_path":"preprints/The-p-adic-section-conjecture-September-24-2026/README.md","date":"2026-09-24","theorem_ref":"thm:section (sections/01-introduction.tex:31); cor:global-modular-curves (sections/08-global-sections.tex:212)","statement_sources":[{"path":"preprints/The-p-adic-section-conjecture-September-24-2026/build/sections/01-introduction.tex","line":31,"label":"thm:section"},{"path":"preprints/The-p-adic-section-conjecture-September-24-2026/build/sections/08-global-sections.tex","line":212,"label":"cor:global-modular-curves"}]}],"lean":null,"reasoning_trace":null}]}
\documentclass[11pt]{article}
\usepackage[margin=25mm]{geometry}
\usepackage{fontspec}
\setmainfont{Latin Modern Roman}
\usepackage{amsmath,amssymb,mathtools}
\usepackage{unicode-math}
\setmathfont{Latin Modern Math}
\usepackage{enumitem}
\usepackage{xurl,hyperref,bookmark}
\hypersetup{colorlinks=true,urlcolor=blue}
\setcounter{secnumdepth}{0}
\setlength{\parindent}{0pt}
\setlength{\parskip}{5pt}
\setlength{\emergencystretch}{3em}
\begin{document}
\section*{The Generalized Étale-Homotopy Section Conjecture}

\subsection{OpenAI's claim}
OpenAI presents the following claims relevant to this question, with manuscript numbers following paper order in Sources. In manuscript 1 (family 019; 2026-09-24), the asserted result is that sections of proper hyperbolic curves over finite extensions of $\mathbb Q_p$ come from points, with a conditional global application through finite-cover descent.

\subsection{Mathematical statement of the claimed result}
\paragraph{Manuscript 1.}
For every prime $p$, finite extension $k/\mathbb Q_p$ take a proper smooth curve $X/k$ that remains connected over $\bar k$ and has genus $g(X)\ge2$. consider the exact sequence of profinite groups
\[
 1\longrightarrow\pi_1(X_{\bar k})\longrightarrow\pi_1(X)
 \longrightarrow G_k\longrightarrow1.
\]
The claimed point-to-section map is a bijection
\[
 X(k)\ \longrightarrow\
 \{s:G_k\to\pi_1(X):s\text{ continuous homomorphism},\;
             (\pi_1(X)\to G_k)s=\operatorname{id}\}/\pi_1(X_{\bar k}),
\]
where sections are identified by conjugation by the full geometric group. The cited modular-curve corollary asserts the same bijectivity over $k=\mathbb Q$ for each smooth projective $X_0(N)$ or $X_1(N)$ of genus at least two, for every positive integer $N$ satisfying that genus condition.

\subsection{Outline of the argument}
For manuscript 1 (family 019; 2026-09-24), two stages of auxiliary covers construct roots of arbitrary order on collars while controlling horizontal ramification. Skeletal cycles and Jacobian torus characters complete the universal covering tower, after which the completion criterion identifies point sections.

\subsection{Scope relative to this problem}
The local theorem concerns proper curves of genus at least two over finite extensions of $\mathbb Q_p$. Global conclusions need the stated finite-descent condition, with the listed modular-curve applications. The unrestricted number-field, nonproper, or general homotopy-section clauses in this question go beyond those cases.

\subsection{Formal verification}
No Lean formalization is listed at this snapshot for manuscript 1 (family 019; 2026-09-24).

This site has not executed the Lean code. Described statement coverage is not an independent proof check.

\subsection{Status caveat}
The manuscripts were produced by an internal OpenAI model. OpenAI cautions that results without formalization can contain mistakes. This site has not independently checked the proofs; the label records the scope of OpenAI's claim.

\subsection{Sources}
\begin{itemize}
\item OpenAI, \emph{The p-adic section conjecture}, 2026-09-24; citation key \texttt{OAI:The-p-adic-section-conjecture-September-24-2026}. \href{https://github.com/openai/math/blob/adc7f1241b42e322a6451854ab7e4b4c146bf78a/preprints/The-p-adic-section-conjecture-September-24-2026/main.pdf}{PDF (GitHub)}; \href{https://github.com/openai/math/raw/adc7f1241b42e322a6451854ab7e4b4c146bf78a/preprints/The-p-adic-section-conjecture-September-24-2026/main.pdf}{PDF (direct)}; \href{https://github.com/openai/math/blob/adc7f1241b42e322a6451854ab7e4b4c146bf78a/preprints/The-p-adic-section-conjecture-September-24-2026/README.md}{Manuscript page and citation}; \href{https://github.com/openai/math/blob/main/preprints/The-p-adic-section-conjecture-September-24-2026/main.pdf}{Latest version}; Statement source: \href{https://github.com/openai/math/blob/adc7f1241b42e322a6451854ab7e4b4c146bf78a/preprints/The-p-adic-section-conjecture-September-24-2026/build/sections/01-introduction.tex\#L31}{thm:section}; Statement source: \href{https://github.com/openai/math/blob/adc7f1241b42e322a6451854ab7e4b4c146bf78a/preprints/The-p-adic-section-conjecture-September-24-2026/build/sections/08-global-sections.tex\#L212}{cor:global-modular-curves}.
\item \href{https://github.com/openai/math/blob/adc7f1241b42e322a6451854ab7e4b4c146bf78a/CONTENTS.md}{OpenAI catalogue entry}.
\end{itemize}
\end{document}
